% ============================================================================
% RASCUNHO COMENTADO: Introdução
% Para publicar, remova o "%" do início de cada linha.
% ============================================================================
%
\chapter{Introdução}
\label{cap:introducao}
%
% A geração fotovoltaica conectada à rede depende de medições contínuas para
% acompanhar a produção, detectar falhas e gerir a usina. Em uma usina
% monitorada, cada inversor reporta a potência ativa entregue à rede, e essa
% série alimenta os indicadores de desempenho. Quando a telemetria falha, o
% registro desaparece, e a lacuna compromete a contabilidade da geração, o
% cálculo de indicadores e a identificação de anomalias.
%
% O problema aparece nos dados da Rede Solarimétrica de Pernambuco
% (RedeSolarPE), mantida pelo Instituto Federal de Pernambuco (IFPE),
% \textit{campus} Pesqueira, que consolida registros operacionais de usinas
% fotovoltaicas conectadas à rede no estado
% \cite{pedrosafilho2022,pedrosafilho2025}. Nos três inversores de 2~kW
% analisados neste trabalho, E1, E2 e E3, a fração de períodos diurnos sem
% registro fica em torno de 13,5\%, e a maior lacuna contínua chega a 12~h.
%
% A proposta deste trabalho é estimar a potência CA do inversor a partir de
% variáveis meteorológicas medidas de forma independente, que continuam
% disponíveis durante a falha. O modelo funciona como um soft-sensor: recebe
% a irradiância transposta para o plano dos módulos, a temperatura ambiente e
% a temperatura do módulo, e devolve a potência estimada no mesmo instante. A
% formulação é a de estimação contemporânea, sem variáveis defasadas
% \cite{moraesneto2024}: a lacuna é preenchida no momento em que ocorre.
%
% O modelo embarcado decorre do cenário de aplicação: a estimativa precisa
% acontecer no local, junto ao inversor, sem depender de um servidor remoto
% e sem interromper a operação quando a rede de comunicação também falha. O
% ESP32-S3 \cite{espressif2024} é o alvo, e o firmware é compilado e gravado
% com o PlatformIO \cite{platformio2024}.
%
% Dois modelos são comparados: o Random Forest (RF), que representa a família
% de árvores e é simples de embarcar como código C \cite{breiman2001}, e o
% perceptron de múltiplas camadas (MLP), que representa as redes neurais e
% pode ser quantizado para inteiros de 8 bits \cite{tflitemicro2024}. A
% comparação é feita sobre o mesmo conjunto, o mesmo \textit{split} e as
% mesmas métricas, em quatro granularidades: 5~min, 10~min, 1~h e 1~d. A
% literatura mostra que árvores e redes neurais têm comportamentos distintos
% em dados tabulares \cite{grinsztajn2022}, e o trabalho verifica qual das
% duas famílias se ajusta melhor ao problema e ao alvo embarcado.
%
% \section{Objetivos}
%
% O objetivo geral é desenvolver e avaliar um soft-sensor embarcado capaz de
% estimar a potência CA de inversores fotovoltaicos a partir de variáveis
% meteorológicas, cobrindo falhas de telemetria.
%
% Os objetivos específicos são:
%
% \begin{alineas}
%   \item preparar o conjunto de dados em quatro granularidades temporais, com
%   separação temporal entre treino, calibração e teste;
%   \item treinar e comparar o Random Forest e o MLP nas quatro
%   granularidades, com seleção de configuração na janela de calibração;
%   \item gerar as bibliotecas embarcadas dos dois modelos e verificar a
%   equivalência entre a implementação em C e o modelo de origem;
%   \item avaliar os modelos no computador e no ESP32-S3, com e sem
%   \textit{gaps} artificiais de telemetria;
%   \item analisar a transferência dos modelos entre granularidades e
%   comparar o erro dentro dos \textit{gaps} com baselines de imputação.
% \end{alineas}
%
% \section{Contribuições}
%
% O trabalho entrega um pipeline completo, do dado bruto ao ensaio na placa,
% com quatro contribuições principais. A primeira é a comparação controlada
% entre um modelo baseado em árvores e uma rede neural, nas mesmas condições
% de dados e avaliação. A segunda é a análise cruzada de granularidades, que
% mostra se um modelo treinado em uma resolução pode ser aplicado a outra. A
% terceira é o ensaio na placa com \textit{gaps} artificiais, que mede o
% desempenho no cenário que justifica o soft-sensor e compara o resultado com
% a interpolação linear e a persistência. A quarta é a automação da geração
% dos artefatos, que reduz o risco de embarcar uma biblioteca desatualizada.
%
% \section{Organização do Trabalho}
%
% O capítulo 2 apresenta a fundamentação teórica: as variáveis ambientais da
% geração fotovoltaica, o conceito de soft-sensor, os dois modelos, as
% métricas de avaliação e as tecnologias de implantação embarcada. O capítulo
% 3 descreve a metodologia: a base de dados, a análise exploratória, a
% preparação, o treino, a avaliação e o ensaio. O capítulo 4 apresenta os
% resultados offline, a análise cruzada e o ensaio na placa. O capítulo 5
% conclui o trabalho, com as limitações e os desdobramentos possíveis.
